% SPDX-License-Identifier: MPL-2.0   (see LICENSE in the AdvancedLatexTemplate repository)
% ============================================================================
% atl-macros.tex -- text helpers and a small, general engineering/maths notation set
%
% Conventions (IEEE style): variables italic; descriptive subscripts, operators and
% units upright; vectors and matrices bold italic through ONE command each (\vect,
% \mat) so the look can be changed in one place. Units always go through siunitx:
%   \valunit{10}{\kilo\ohm}   \qty{2.2}{\micro\farad}   \num{1.5e-3}
% No LaTeX built-ins are redefined (\H \P \Re \Im \vec \div ...).
% Domain-specific packs (electronics, power, RF, ML ...) belong in preamble/local.tex.
% ============================================================================
\makeatletter
% ---- text helpers ----------------------------------------------------------------
\newcommand{\keyterm}[1]{\textbf{\textcolor{atl-primary}{#1}}}
\newcommand{\alert}[1]{\textbf{\textcolor{atl-warning}{#1}}}
\newcommand{\ok}[1]{\textbf{\textcolor{atl-success}{#1}}}
\newcommand{\highlight}[1]{\colorbox{atl-accent!28}{\strut #1}}
\newcommand{\muted}[1]{\textcolor{atl-muted}{#1}}
\newcommand{\todo}[1]{\textcolor{atl-warning}{\textbf{[TODO: #1]}}}

% ---- notation ---------------------------------------------------------------------
\newcommand{\qtag}[2]{\ensuremath{#1_{\mathrm{#2}}}}              % R_in
\newcommand{\qindex}[2]{\ensuremath{#1_{#2}}}                       % x_n
\newcommand{\qtagindex}[3]{\ensuremath{#1_{\mathrm{#2},#3}}}       % I_ref,k
\newcommand{\vect}[1]{\ensuremath{\bm{#1}}}
\newcommand{\mat}[1]{\ensuremath{\bm{#1}}}
\newcommand{\uvec}[1]{\ensuremath{\hat{\bm{#1}}}}
\newcommand{\phasor}[1]{\ensuremath{\tilde{#1}}}
\newcommand{\conj}[1]{\ensuremath{\overline{#1}}}
\newcommand{\trans}[1]{\ensuremath{{#1}^{\mathsf{T}}}}
\newcommand{\herm}[1]{\ensuremath{{#1}^{\mathsf{H}}}}
\newcommand{\pinv}[1]{\ensuremath{{#1}^{\dagger}}}
\DeclarePairedDelimiter{\abs}{\lvert}{\rvert}
\DeclarePairedDelimiter{\norm}{\lVert}{\rVert}
\DeclarePairedDelimiter{\inner}{\langle}{\rangle}
\newcommand{\ddiff}{\ensuremath{\mathop{}\!\mathrm{d}}}
\newcommand{\deriv}[2]{\ensuremath{\frac{\mathrm{d}#1}{\mathrm{d}#2}}}
\newcommand{\nderiv}[3]{\ensuremath{\frac{\mathrm{d}^{#3}#1}{\mathrm{d}#2^{#3}}}}
\newcommand{\pderiv}[2]{\ensuremath{\frac{\partial#1}{\partial#2}}}
\newcommand{\integral}[4]{\ensuremath{\int_{#1}^{#2}#3\,\mathrm{d}#4}}
\newcommand{\econst}{\ensuremath{\mathrm{e}}}
\newcommand{\junit}{\ensuremath{\mathrm{j}}}
\newcommand{\iunit}{\ensuremath{\mathrm{i}}}
\newcommand{\jw}{\ensuremath{\junit\omega}}
\newcommand{\N}{\ensuremath{\mathbb{N}}}
\newcommand{\Z}{\ensuremath{\mathbb{Z}}}
\newcommand{\Q}{\ensuremath{\mathbb{Q}}}
\newcommand{\R}{\ensuremath{\mathbb{R}}}
\newcommand{\C}{\ensuremath{\mathbb{C}}}
\DeclareMathOperator{\sgn}{sgn}
\DeclareMathOperator{\tr}{tr}
\DeclareMathOperator{\rank}{rank}
\DeclareMathOperator{\diag}{diag}
\DeclareMathOperator{\sinc}{sinc}
\DeclareMathOperator{\rect}{rect}
\DeclareMathOperator{\Expect}{\mathbb{E}}
\DeclareMathOperator{\Var}{Var}
\DeclareMathOperator{\Cov}{Cov}
\DeclareMathOperator{\Prob}{\mathbb{P}}
\DeclareMathOperator*{\argmin}{arg\,min}
\DeclareMathOperator*{\argmax}{arg\,max}

% ---- units (siunitx v3) -------------------------------------------------------------
\newcommand{\valunit}[3][]{\qty[#1]{#2}{#3}}
\DeclareSIUnit{\dBm}{dBm}
\DeclareSIUnit{\var}{var}
\DeclareSIUnit{\rpm}{rpm}
\DeclareSIUnit{\ppm}{ppm}
\DeclareSIUnit{\pu}{p.u.}
\makeatother
